\textbf{Phần 1. Trắc nghiệm lựa chọn}
\setcounterpageref{ex}{0}

\begin{ex} %[3H5N34]
Để đo chiều dài của một con kiến, ta nên dùng đơn vị nào?
\choice
{Mét ($\text{m}$)}
{Đề-xi-mét ($\text{dm}$)}
{\True Mi-li-mét ($\text{mm}$)}
{Ki-lô-mét ($\text{km}$)}
\loigiai{
Đo vật nhỏ như con kiến dùng đơn vị $\text{mm}$.\\
\textbf{Đáp án của câu hỏi là C}
}
\end{ex}

\begin{ex} %[3H5N34]
Để đo lượng nước trong một ống tiêm, ta nên dùng đơn vị nào?
\choice
{\True Mi-li-lít ($\text{ml}$)}
{Lít ($\text{l}$)}
{Gam ($\text{g}$)}
{Mét ($\text{m}$)}
\loigiai{
Đo lượng dung dịch nhỏ dùng mi-li-lít ($\text{ml}$).\\
\textbf{Đáp án của câu hỏi là A}
}
\end{ex}

\begin{ex} %[3H5H34]
Một quả trứng gà nặng khoảng
\choice
{$50\text{ kg}$}
{$500\text{ g}$}
{$5\text{ g}$}
{\True $50\text{ g}$}
\loigiai{
Quả trứng gà nặng khoảng $50\text{ g}$.\\
\textbf{Đáp án của câu hỏi là D}
}
\end{ex}

\begin{ex} %[3H5H34]
Bề dày của một chiếc thẻ căn cước (ATM) dày khoảng
\choice
{$1\text{ cm}$}
{\True $1\text{ mm}$}
{$10\text{ mm}$}
{$5\text{ cm}$}
\loigiai{
Thẻ ATM dày khoảng $1\text{ mm}$.\\
\textbf{Đáp án của câu hỏi là B}
}
\end{ex}

\begin{ex} %[3H5H34]
Một hộp sữa chua thường có khối lượng khoảng
\choice
{$10\text{ g}$}
{$1\text{ kg}$}
{$1000\text{ g}$}
{\True $100\text{ g}$}
\loigiai{
Hộp sữa chua khoảng $100\text{ g}$.\\
\textbf{Đáp án của câu hỏi là D}
}
\end{ex}

\begin{ex} %[3H5H34]
Tính: $250\text{ g} \times 3 + 150\text{ g} = \dots$
\choice
{\True $900\text{ g}$}
{$800\text{ g}$}
{$1000\text{ g}$}
{$850\text{ g}$}
\loigiai{
$750 + 150 = 900\text{ g}$.\\
\textbf{Đáp án của câu hỏi là A}
}
\end{ex}

\begin{ex} %[3H5V34]
Một bình chứa $1\text{ l}$ nước, người ta rót vào bình thêm $350\text{ ml}$ si-rô. Tổng lượng dung dịch trong bình là
\choice
{$1035\text{ ml}$}
{\True $1350\text{ ml}$}
{$1300\text{ ml}$}
{$450\text{ ml}$}
\loigiai{
$1000 + 350 = 1350\text{ ml}$.\\
\textbf{Đáp án của câu hỏi là B}
}
\end{ex}

\begin{ex} %[3H5V34]
Một thanh kim loại dài $12\text{ cm } 5\text{ mm}$. Sau khi nung nóng nó dài thêm $2\text{ mm}$. Chiều dài lúc này là
\choice
{\True $12\text{ cm } 7\text{ mm}$}
{$12\text{ cm } 3\text{ mm}$}
{$14\text{ cm } 5\text{ mm}$}
{$127\text{ cm}$}
\loigiai{
$12\text{ cm } 5\text{ mm} + 2\text{ mm} = 12\text{ cm } 7\text{ mm} = 127\text{ mm}$.\\
\textbf{Đáp án của câu hỏi là A}
}
\end{ex}

\begin{ex} %[3H5V34]
Mẹ mua 4 gói mì tôm, mỗi gói $75\text{ g}$. Tổng khối lượng 4 gói mì tôm là
\choice
{$280\text{ g}$}
{$320\text{ g}$}
{\True $300\text{ g}$}
{$350\text{ g}$}
\loigiai{
$75 \times 4 = 300\text{ g}$.\\
\textbf{Đáp án của câu hỏi là C}
}
\end{ex}

\begin{ex} %[3H5V34]
Một chai dầu gội đầu chứa $650\text{ ml}$. Sau 1 tháng sử dụng còn lại $200\text{ ml}$. Lượng dầu gội đã dùng là
\choice
{$400\text{ ml}$}
{\True $450\text{ ml}$}
{$500\text{ ml}$}
{$420\text{ ml}$}
\loigiai{
$650 - 200 = 450\text{ ml}$.\\
\textbf{Đáp án của câu hỏi là B}
}
\end{ex}

\begin{ex} %[3H5V34]
Nhiệt độ ngoài sân lúc 12 giờ trưa là $36^\circ\text{C}$, trong phòng bật quạt là $31^\circ\text{C}$. Trong phòng mát hơn ngoài sân bao nhiêu độ C?
\choice
{$4^\circ\text{C}$}
{$6^\circ\text{C}$}
{$7^\circ\text{C}$}
{\True $5^\circ\text{C}$}
\loigiai{
$36 - 31 = 5^\circ\text{C}$.\\
\textbf{Đáp án của câu hỏi là D}
}
\end{ex}

\begin{ex} %[3H5C34]
Một chai nước ngọt nặng tất cả $1\text{ kg } 100\text{ g}$ khi đầy nước. Riêng vỏ chai rỗng nặng $100\text{ g}$. Biết $1\text{ ml}$ nước ngọt nặng đúng $1\text{ g}$. Dung tích nước ngọt trong chai là
\choice
{$900\text{ ml}$}
{$1100\text{ ml}$}
{\True $1000\text{ ml}$ (hay $1\text{ l}$)}
{$800\text{ ml}$}
\loigiai{
Lượng nước ngọt nặng: $1100 - 100 = 1000\text{ g} = 1000\text{ ml}$.\\
\textbf{Đáp án của câu hỏi là C}
}
\end{ex}

\textbf{Phần 2. Trắc nghiệm đúng sai}
\setcounterpageref{ex}{0}

\begin{ex} %[3H5H34]
Xét tính đúng sai về ước lượng và đo lường thực tế:
\choiceTF
{\True Chiếc bút chì dài khoảng $15\text{ cm}$.}
{\True Một chai nước khoáng nhỏ chứa khoảng $500\text{ ml}$.}
{Cặp sách của học sinh nặng khoảng $30\text{ kg}$.}
{\True Cặp sách của học sinh nặng khoảng $3\text{ kg}$ ($3000\text{ g}$).}
\loigiai{
\begin{itemize}
    \item Bút chì $15\text{ cm}$: Đúng.
    \item Chai nước $500\text{ ml}$: Đúng.
    \item Cặp sách $30\text{ kg}$ là quá nặng (khoảng $3\text{ kg}$). Ý 3 Sai, Ý 4 Đúng.
\end{itemize}
\begin{itemize}
    \item a) \textbf{Đúng}.
    \item b) \textbf{Đúng}.
    \item c) \textbf{Sai}.
    \item d) \textbf{Đúng}.
\end{itemize}
}
\end{ex}

\begin{ex} %[3H5V34]
Trong một thí nghiệm, bạn Lan pha $200\text{ ml}$ nước nóng $80^\circ\text{C}$ với $200\text{ ml}$ nước lạnh $20^\circ\text{C}$. Xét tính đúng sai:
\choiceTF
{\True Tổng dung tích nước thu được là $400\text{ ml}$.}
{\True Nhiệt độ của nước sau khi pha nằm trong khoảng từ $20^\circ\text{C}$ đến $80^\circ\text{C}$.}
{Nhiệt độ nước sau khi pha cao hơn $80^\circ\text{C}$.}
{\True Nhiệt độ nước sau khi pha vào khoảng $50^\circ\text{C}$ (ấm vừa).}
\loigiai{
\begin{itemize}
    \item $200 + 200 = 400\text{ ml}$. Đúng.
    \item Nhiệt độ cân bằng nằm giữa $20$ và $80^\circ\text{C}$. Đúng.
    \item Không thể cao hơn nhiệt độ nước nóng ban đầu. Sai.
    \item Trung bình cộng khoảng $50^\circ\text{C}$. Đúng.
\end{itemize}
\begin{itemize}
    \item a) \textbf{Đúng}.
    \item b) \textbf{Đúng}.
    \item c) \textbf{Sai}.
    \item d) \textbf{Đúng}.
\end{itemize}
}
\end{ex}

\textbf{Phần 3. Trả lời ngắn}
\setcounterpageref{ex}{0}

\begin{ex} %[3H5H34]
Đổi: $8\text{ cm } 4\text{ mm}$ ra mi-li-mét.
\shortans{84}
\loigiai{
\textbf{Đáp số:} $84$\\
$80 + 4 = 84\text{ mm}$.
}
\end{ex}

\begin{ex} %[3H5H34]
Tính: $350\text{ g} + 650\text{ g}$.
\shortans{1000}
\loigiai{
\textbf{Đáp số:} $1000$\\
$350 + 650 = 1000\text{ g}$.
}
\end{ex}

\begin{ex} %[3H5V34]
Tính: $250\text{ ml} \times 4$.
\shortans{1000}
\loigiai{
\textbf{Đáp số:} $1000$\\
$250 \times 4 = 1000\text{ ml}$.
}
\end{ex}

\begin{ex} %[3H5V34]
Nhiệt độ lúc sáng là $21^\circ\text{C}$, buổi trưa là $32^\circ\text{C}$. Hỏi nhiệt độ đã tăng thêm bao nhiêu độ C?
\shortans{11}
\loigiai{
\textbf{Đáp số:} $11$\\
$32 - 21 = 11^\circ\text{C}$.
}
\end{ex}

\begin{ex} %[3H5V34]
Có 3 gói kẹo, mỗi gói nặng $200\text{ g}$ và 1 gói bánh nặng $300\text{ g}$. Tổng khối lượng là bao nhiêu gam?
\shortans{900}
\loigiai{
\textbf{Đáp số:} $900$\\
$200 \times 3 + 300 = 900\text{ g}$.
}
\end{ex}

\begin{ex} %[3H5C34]
Một ca nước chứa $1\text{ l}$ nước. Rót vào 5 cốc như nhau thì còn thừa $150\text{ ml}$. Hỏi mỗi cốc đựng bao nhiêu mi-li-lít nước?
\shortans{170}
\loigiai{
\textbf{Đáp số:} $170$\\
Đã rót: $1000 - 150 = 850\text{ ml}$.\\
Mỗi cốc: $850 : 5 = 170\text{ ml}$.
}
\end{ex}
